
\section{Context-Bound Envelope Defense}
\label{app:defense-proposal}

\looseness=-1 The attacks presented in this paper: cross-session,
cross-user, and cross-model replay, all exploit the same structural gap:
the AEAD envelope authenticates the \emph{content} of a reasoning block
but not the \emph{context} in which it was produced or is later replayed.
We therefore propose a defense that re-binds every envelope to its
originating user, session, and conversational position, complemented by
operational measures for legacy data, backwards compatibility, and
training-time hardening. Table~\ref{tab:defense-summary} maps each leakage vector to exactly
what the corresponding mitigation provides. The remainder of
this section elaborates each entry in turn.

\begin{table}[H]
\centering
\footnotesize
\renewcommand{\arraystretch}{1.35}
\setlength{\tabcolsep}{4pt}
\caption{Leakage vectors mapped to protections and processes.
Each tick states what the mitigation supplies; the final
column records the operational steps that deliver it.}
\label{tab:defense-summary}
\begin{tabular}{@{}
  >{\centering\arraybackslash}p{0.035\textwidth}
  >{\raggedright\arraybackslash}p{0.15\textwidth}
  >{\raggedright\arraybackslash}p{0.34\textwidth}
  >{\raggedright\arraybackslash}p{0.38\textwidth}
  @{}
}
\toprule
\textbf{\#} & \textbf{Issue} & \textbf{What the mitigation provides} & \textbf{Mitigation} \\
\midrule
1 &
Cross-user leakage &
$\checkmark$ User-identity binding \newline
$\checkmark$ Stateless verification \newline
$\checkmark$ Immediate mismatch rejection
&
\begin{enumerate}[topsep=0pt,partopsep=0pt,before=\vspace{-0.55\baselineskip}]
  \item Embed \texttt{user\_id} in AEAD associated data at issuance.
  \item On replay, compare bound identity to authenticated caller.
  \item Reject the envelope on any mismatch.
\end{enumerate} \\
\addlinespace
2 &
Cross-session leakage &
$\checkmark$ Session + predecessor binding \newline
$\checkmark$ Ordinality (P1) under compaction \newline
$\checkmark$ Native fork / compact / downgrade support \newline
$\checkmark$ Dramatically reduced blast radius
&
\begin{enumerate}[topsep=0pt,partopsep=0pt,before=\vspace{-0.55\baselineskip}]
  \item Hash-chain each envelope to \texttt{session\_id} and its
        predecessor (Eq.~\ref{eq:chain}).
  \item Enforce ordinality server-side.
  \item Retain only Merkle roots after compaction so surviving
        spans stay verifiable.
\end{enumerate} \\
\addlinespace
3 &
Legacy public/enterprise datasets &
$\checkmark$ Permanent undecodability of pre-fix signatures \newline
$\checkmark$ Clean cryptographic separation from new material
&
\begin{enumerate}[topsep=0pt,partopsep=0pt,before=\vspace{-0.55\baselineskip}]
  \item Rotate every pre-fix signing key.
  \item Refuse to decode any envelope under a retired key ID.
  \item (Optional) Offer identity-verified re-signing for enterprise
        archives.
\end{enumerate} \\
\addlinespace
4 &
Backwards compatibility &
$\checkmark$ Zero-break migration path \newline
$\checkmark$ Bounded dual-format window \newline
$\checkmark$ Identity-verified re-issuance
&
\begin{enumerate}[topsep=0pt,partopsep=0pt,before=\vspace{-0.55\baselineskip}]
  \item Accept both legacy and context-bound envelopes during a
        fixed deprecation window.
  \item Expose an opt-in batch re-signature endpoint.
  \item Re-issue only after confirming the requester owns the
        original session.
\end{enumerate} \\
\addlinespace
5 &
Model-level compliance &
$\checkmark$ Closure of residual gaps beyond cryptography \newline
$\checkmark$ Resistance to transcription / replay jailbreaks
&
\begin{enumerate}[topsep=0pt,partopsep=0pt,before=\vspace{-0.55\baselineskip}]
  \item Post-train models to recognise transcription-style prompts
        (e.g., \texttt{<thinking-copy>}).
  \item Refuse the request irrespective of envelope validity.
\end{enumerate} \\
\addlinespace
6 &
Nonce predictability &
$\checkmark$ Critical since key is shared between users \newline
$\checkmark$ Cryptographic foundation for every binding above \newline
$\checkmark$ Collision- and forgery-resistance at provider scale 
&
\begin{enumerate}[topsep=0pt,partopsep=0pt,before=\vspace{-0.55\baselineskip}]
  \item Draw a high-entropy nonce from a CSPRNG for every block.
  \item Enforce server-side uniqueness before the envelope is issued.
\end{enumerate} \\
\bottomrule
\end{tabular}
\end{table}

\subsection{Cross-User Binding}
\label{sec:cross-user-binding}

The simplest and cheapest fix addresses Issue~1 directly: embed the
originating \texttt{user\_id} (or an equivalent stable account
identifier) inside the AEAD associated data of every reasoning envelope.
On replay the API simply compares the bound identifier against the
already-authenticated caller and rejects the block on mismatch.
This closes the cross-user attack vector entirely without requiring any stateful backend.

\subsection{Chained, Context-Bound Envelopes}
\label{sec:chained-binding}

Cross-session replay (Issue~2) is harder, because legitimate workflows
depend on the very portability that enables the attack: users must be
able to \emph{fork} a conversation, \emph{compact} old turns out of a
long session, and \emph{downgrade} to a cheaper model mid-conversation.
A naive binding of every block to the literal, complete prior transcript
would break all three.

We therefore propose a lightweight hash chain that binds each user turn and reasoning
block only to its session and its immediate predecessor:
\begin{equation}
\label{eq:chain}
\tau_{n+1}
\;=\;
H\bigl(\,
  \text{user\_id}
  \,\|\,
  \text{session\_id}
  \,\|\,
  H(\tau_n \,\|\, \text{salt}_2)
  \,\|\,
  \text{salt}_1
\,\bigr),
\end{equation}
where $\tau_n$ denotes the reasoning content of block~$n$ and the resulting
hash is placed in the associated data of the following block's envelope.

Chaining alone does not prevent an adversary who replays an \emph{entire}
prior conversation, in order, from coercing a compatible decoder into
transcription.  What it does achieve is a substantial increase in attack
cost (the adversary now needs the full session rather than a single
captured signature) together with a far smaller, auditable blast radius---
sufficient to defeat the scalable, single-signature extraction attacks.

\paragraph{Design properties.}
A practical chained scheme could satisfy four complementary properties:

\begin{description}
\item[P1 (Reasoning Ordinality).]
  Block~$X$ must be verifiably prior to block~$Y$ whenever~$Y$ depends
  on~$X$.  This weaker ordering guarantee is inexpensive to preserve under
  compaction: dropped chunks can simply be excised from the chain while
  the relative order of surviving nodes remains intact.

\item[P2 (Full Reasoning Integrity).]
  Block~$X$ must be preceded by \emph{exactly} block~$X{-}1$; no
  intermediate block may be elided.  The stronger guarantee is costly
  under compaction, because removal of an internal node forces
  recomputation of every downstream hash.

\item[P3 (Leak monitoring).]
  Providers should run a continuous verbatim-matching pass between any
  published or reported plaintext and any decoded reasoning that surfaces
  through support channels or abuse reports, so that leakage is detected and account flagged rather than merely prevented.

\item[P4 (Non-replayability).]
  The same reasoning block must never be accepted twice -- neither within
  a session nor across sessions.  Replay of an already-consumed envelope
  is rejected server-side.
\end{description}

Properties P1 and P2 stand in tension with the compaction requirement;
we therefore do not attempt to guarantee P2 globally.  Instead we recommend
a Merkle-tree structure over the session's blocks: the leaves are the
individual reasoning envelopes, while the provider retains only the tree
root (or a small number of subtree roots) after leaves are pruned.
This yields P1 cheaply everywhere and P2 on demand for any surviving
contiguous span, without ever requiring the full unpruned history to be
replayed.

Model downgrade and forking are handled uniformly by allowing a session
to branch the tree: a fork inherits the chain state up to the fork point
and thereafter continues independently. Consequently an envelope's
context binding never needs to reference a model version that produced
an earlier segment of an unrelated branch.

\subsection{Legacy Data}
\label{sec:legacy}

None of the bindings above can protect data that has already been
produced and published, e.g. the 6\,708 public trajectories
surveyed.  Because old envelopes
were signed under a key that encodes neither user nor session context,
the only retroactive remedy is key invalidation: providers must rotate
all pre-fix signing keys and refuse to decode any envelope authenticated
under a retired key ID.  This single measure -- independent of any new
context-binding logic -- renders previously published signatures
permanently undecodable, at the inevitable cost of also invalidating
legitimate continuations of old sessions.

\subsection{Backwards Compatibility}
\label{sec:backcompat}

Enterprise customers hold large volumes of stored transcripts whose
signatures they may still need to resume (e.g., paused agentic workflows).
We therefore recommend a bounded dual-format acceptance window during
which both legacy and context-bound envelopes are honored, together with
an opt-in batch re-signature endpoint.  Customers may submit archived
transcripts for re-issuance under the new scheme, subject to identity
verification that the requesting account is the original session owner.
Once the deprecation window closes, legacy envelopes are rejected
outright, consistent with the key-rotation policy of
Section~\ref{sec:legacy}.

\subsection{Training-Time Defenses}
\label{sec:training}

Cryptographic binding constrains \emph{which} model may be asked to
decode a given envelope; it cannot constrain what a compliant decoder
does once it is legitimately asked to process its own prior reasoning. 
Closing this residual gap is therefore a training problem. Models should be explicitly trained to recognize and
refuse transcription-style jailbreaks (e.g., \texttt{<thinking-copy>}
framings) regardless of how innocuous the surrounding request appears.
We list this as a standing item for future post-training work rather than
a solved component of the present proposal.